\pagestyle{plain}
\chapter{Introdução}

\section{Saberes e Sistemas Agrícolas Tradicionais}

Os Saberes e Sistemas Agrícolas Tradicionais (SSAT) não se reduzem a práticas agrárias pretéritas ou à agricultura de subsistência. Sob a ótica da agroecologia e da etnoecologia, esses sistemas constituem complexos ecossistêmicos antropogênicos, forjados pela co-evolução milenar entre grupos humanos e seus biomas. Toledo e Barrera-Bassols \cite{Toledo2008} conceituam essa interação como a expressão da memória biocultural, uma matriz indissociável onde a diversidade biológica é mantida e amplificada pela diversidade cultural, linguística e simbólica das comunidades. Nesse contexto, o saber não opera como um arquivo estático, mas estrutura-se como um \textit{corpus} dinâmico fundamentado no complexo Conhecimento-Prática-Crença (K-P-B) descrito por Fikret Berkes \cite{Berkes2017}.

Sob esta ótica analítica, o componente de crenças, frequentemente marginalizado por análises tecnocráticas como elemento subjetivo, desempenha uma função estruturante crítica. Ele atua como uma cosmovisão normativa que impõe restrições éticas à apropriação dos recursos naturais, transformando a paisagem em um sujeito de direitos e reciprocidade, e não mero objeto de exploração \cite{Toledo2008}. A indissociabilidade dessa tríade engendra um sistema cibernético de retroalimentação social e ecológica, onde a observação empírica refinada (\textit{Knowledge}) instrumentaliza as instituições de manejo (\textit{Practice}), as quais são, por sua vez, moduladas e legitimadas por valores culturais (\textit{Belief}) \cite{Berkes2017,Ostrom1990}. Esse arranjo sistêmico impede a sobre-exploração dos estoques biológicos através de tabus e sanções sociais, assegurando que a inovação técnica permaneça subordinada à integridade dos ciclos ecossistêmicos que sustentam a comunidade \cite{Berkes2003}.

A racionalidade produtiva que rege esses sistemas difere radicalmente da lógica industrial de simplificação ecológica. Enquanto a agricultura moderna busca a homogeneização e o isolamento de variáveis, os SSAT operam sob a lógica da complexidade, da diversidade funcional e da redundância biológica. Miguel Altieri \cite{Altieri1995} descreve esses agroecossistemas como exemplos de engenharia ecológica vernácula, caracterizados por arranjos policulturais, sistemas agroflorestais e manejo integrado de paisagens que maximizam as interações sinérgicas entre solo, plantas e animais. Essa sofisticação sistêmica permite que tais comunidades mantenham altos níveis de produtividade líquida e estabilidade a longo prazo sem dependência massiva de insumos externos, constituindo verdadeiros refúgios de agrobiodiversidade e bancos genéticos \textit{in situ} vitais para o futuro da agricultura global.

Portanto, compreender esses sistemas exige uma abordagem holística que reconheça a sua multifuncionalidade. Eles não produzem apenas \textit{commodities}, mas geram serviços ecossistêmicos quantificáveis \cite{Costanza1997}, reproduzem identidades territoriais \cite{CalvetMir2016} e sustentam estruturas sociais de cooperação cuja vitalidade condiciona a resiliência socioecológica do agroecossistema \cite{Folke2016}. A ``tradicionalidade'' aqui não denota atraso tecnológico, mas uma trajetória tecnológica distinta, pautada na observação fenológica rigorosa e na experimentação contínua \cite{GomezBaggethun2013}. É justamente essa riqueza de variáveis e interações não-lineares, definida como a complexidade inerente à memória biocultural, que desafia os métodos reducionistas de análise e justifica a adoção de ferramentas computacionais avançadas capazes de decodificar a inteligência ecológica neles depositada \cite{Berkes2003}.

A gestão estratégica de ativos intangíveis consolidou-se como diferencial competitivo na economia do conhecimento, exigindo novas abordagens para a valoração e proteção de saberes ancestrais. A fundamentação teórica desta premissa remonta à teoria do crescimento da firma de Penrose \cite{Penrose1959}, que identificou a base de recursos internos e as capacidades gerenciais como os limitadores primários e, simultaneamente, os motores do crescimento organizacional. Transpondo essa lógica para o contexto comunitário, os saberes tradicionais deixam de ser vistos como meros traços culturais para serem compreendidos como recursos estratégicos vitais. Na perspectiva da Resource-Based View (RBV) \cite{Barney1991}, a vantagem competitiva sustentável e a sobrevivência de um grupo derivam de recursos que sejam Valiosos, Raros, Inimitáveis e Não-substituíveis (VRIN). O conhecimento tradicional, quando adequadamente gerido e protegido contra apropriações indébitas, atende rigorosamente a esses critérios, configurando-se como um ativo idiossincrático de alto valor.

Contudo, a análise desses ativos demanda superar a visão estática de estoque de conhecimento. Sob a lente da teoria econômica evolucionária, os Saberes e Sistemas Agrícolas Tradicionais (SSAT) ilustram o conceito de novas combinações proposto por Joseph Schumpeter \cite{Schumpeter1934}. Longe de representarem o oposto da modernidade, esses sistemas constituem laboratórios vivos de inovação contínua, onde a persistência secular de práticas agroecológicas em ambientes de restrição de recursos evidencia um vigoroso processo de experimentação empírica. Nessa perspectiva analítica, o acervo de conhecimentos detido pelas comunidades opera como um conjunto de Capacidades Dinâmicas, conforme a estrutura teórica de Teece, Pisano e Shuen \cite{Teece1997}. O monitoramento fenológico minucioso corresponde ao mecanismo de \textit{sensing} (monitoramento de sinais fracos), a seleção de genótipos adaptados opera como o \textit{seizing} (mobilização de recursos) e a modificação de práticas de cultivo reflete o \textit{transforming} (reconfiguração de ativos).

O desafio central para a ciência contemporânea reside na natureza epistemológica desses ativos. Conforme postulado por Michael Polanyi \cite{Polanyi1966}, grande parte desse conhecimento reside em uma dimensão tácita e inefável, operando sob a lógica de que ``sabemos mais do que somos capazes de expressar''. Diferentemente da inovação formal, codificada em patentes e manuais, a inovação nos sistemas tradicionais é transmitida oralmente e incorporada na prática cotidiana (\textit{learning-by-doing}), o que frequentemente a torna invisível às métricas da ciência hegemônica e aos instrumentos clássicos de propriedade intelectual. A crise ambiental global e os imperativos do Antropoceno, entretanto, revelam que esses modelos de gestão tácita incorporam dimensões socioambientais de resiliência que a agricultura industrial falhou em replicar, tornando urgente o desenvolvimento de mecanismos de ``externalização'', conforme o modelo de Nonaka \cite{Nonaka1995}, capazes de traduzir essa complexidade para linguagens validáveis.

Operacionalizar essa visão de conhecimentos tradicionais como ativos estratégicos e dinâmicos demanda, portanto, instrumentos jurídicos e metodologias tecnológicas de decodificação. Se o conhecimento tradicional opera como um sistema adaptativo complexo e \textit{path-dependent} \cite{Nelson1982}, sua gestão eficaz exige o desenvolvimento de capacidades absortivas organizacionais \cite{CohenLevinthal1990}, permitindo que as comunidades integrem tecnologias digitais sem comprometer sua autenticidade cultural. O paradigma da Inovação Aberta \cite{Chesbrough2003} fornece o arcabouço para essa integração, onde o fluxo de conhecimento entre a ciência de dados e a sabedoria ancestral não ocorre por imposição, mas por processos de co-criação que visam salvaguardar o patrimônio biocultural.

A valoração de conhecimentos tradicionais como ativos intangíveis de inovação demanda metodologias que superem abordagens econométricas lineares ou puramente baseadas em custos. A literatura sobre valoração de ativos evoluiu para frameworks analíticos que incorporam a complexidade multidimensional da inovação, conforme discutido por Crossan e Apaydin \cite{Crossan2010}. Nesse cenário, a densidade de interações biológicas, climáticas e sociais presentes nos sistemas tradicionais excede a capacidade de apreensão por métodos de observação estatística convencional, o que reforça a pertinência de ferramentas de Engenharia de Dados e Inteligência Artificial.

O Aprendizado de Máquina (Machine Learning) emerge neste contexto não apenas como ferramenta instrumental de processamento, mas como decodificador tecnológico de capacidades dinâmicas. Ao processar a alta dimensionalidade de dados fenológicos e ambientais, algoritmos avançados possuem a capacidade de identificar padrões latentes que constituem os microfundamentos da adaptação tradicional \cite{Mountrakis2011,Spyrou2025}. Contudo, conforme demonstrado pela revisão de escopo conduzida no Capítulo~1, embora a acurácia combinada das aplicações de ML a sistemas agrícolas tradicionais alcance 90,7\% (IC~95\% 89,8--91,5\%), a conformidade FAIR permanece criticamente baixa (escore médio 18,7/100) e a ausência de protocolos de validação espacial compromete a prontidão operacional para aplicações regulatórias \cite{Wilkinson2016,Rudin2019}. Essa constatação reforça a necessidade de, antes de aplicar tecnologias computacionais sofisticadas, dispor de variáveis de entrada estruturadas, validadas e contextualmente pertinentes. A aplicação dessas tecnologias permite explicitar o conhecimento tácito, transformando observações empíricas dispersas em parâmetros científicos auditáveis e validáveis \cite{Nonaka1995}. Esse processo de tradução tecnológica é fundamental para fundamentar a legitimidade científica dos SSAT, criando a base probatória necessária para sua proteção jurídica e valorização econômica.

A salvaguarda dos saberes dos Povos e Comunidades Tradicionais (PCTs) representa, portanto, um desafio estratégico de gestão da inovação e de justiça territorial. No Brasil, apesar dos avanços constitucionais e da regulamentação do acesso ao patrimônio genético e aos conhecimentos tradicionais associados \cite{brasil2015lei13123}, a insegurança jurídica territorial permanece crítica, conforme evidenciado pelo relatório que associa 65\% dos assassinatos em territórios quilombolas à ausência de titulação \cite{conaq2023racismo}. A falta de documentação técnica padronizada que comprove a ocupação ancestral e a gestão sofisticada dos recursos naturais frequentemente fragiliza a posição dessas comunidades em disputas fundiárias e processos de licenciamento \cite{Santilli2009}. Nesse sentido, o mapeamento sistemático e a modelagem computacional dos conhecimentos agroecológicos não servem apenas à curiosidade acadêmica, mas constituem evidências materiais robustas de ocupação e uso sustentável do território.


A relevância desta investigação reside na proposição de um modelo integrado que instrumentaliza a Engenharia de Dados para o empoderamento comunitário e a proteção de ativos estratégicos. A gestão de conhecimentos tradicionais, quando suportada por evidências computacionais robustas, fortalece a governança territorial e a soberania epistemológica das comunidades \cite{FalsBorda1991}. O projeto contribui para a consolidação de capacidades territoriais de inovação social, alinhando-se a modelos de hélice tripla \cite{EtzkowitzLeydesdorff2000} que integram universidades, comunidades e instituições de pesquisa em prol do desenvolvimento sustentável.

A vulnerabilidade biocultural dos SSAT não opera, contudo, como vetor unidimensional. A meta-análise quantitativa conduzida no Capítulo~2 demonstra que a erosão intergeracional, mediada por aculturação (substituição \textit{in situ} de práticas endógenas por padrões da cultura dominante) e emigração de detentores, cuja convergência reduz a proporção de jovens que dominam rotinas agroecológicas entre coortes consecutivas \cite{ReyesGarcia2013,GomezBaggethun2013}, desencadeia efeito cascata sobre a complexidade biocultural dos agroecossistemas, pois a simplificação do repertório de manejo reduz a riqueza de espécies cultivadas e elimina interações solo-planta cuja codificação depende de transmissão oral continuada \cite{Jarvis2008}. Esse processo retroalimenta-se pela degradação da organização social comunitária, acelerando a própria erosão que o originou \cite{Agrawal2002}. A investigação opera, assim, sobre oito dimensões funcionais de vulnerabilidade ($V_{1}$--$V_{8}$, detalhadas na Seção Justificativa), cujas magnitudes de efeito e moderadores contextuais subsidiam a ponderação de variáveis no Índice de Salvaguarda Biocultural (ISB).

Nesse contexto, a presente tese estrutura-se sob delineamento metodológico processual e integrativo. A investigação inicia-se por revisão de escopo PRISMA-ScR e meta-análise quantitativa \cite{Tricco2018b,Page2021}, visando diagnosticar a prontidão operacional de ML em contextos agroecológicos e sintetizar as magnitudes de efeito das dimensões de vulnerabilidade biocultural. Subsequentemente, avança para a adaptação transcultural de instrumentos (WOCAT-SLM) e elicitação Delphi, convertendo saberes tácitos em variáveis linguísticas validadas. A cadeia prossegue com modelagem psicométrica (TRI) e motor de inferência fuzzy para construção do ISB \cite{Zadeh1965}, auditoria computacional e avaliação de aceitação institucional, e culmina no design participativo e validação territorial com comunidades quilombolas do Semiárido Nordeste~II (Bahia).

Esta abordagem integrada visa superar a lacuna entre a riqueza empírica dos saberes tradicionais e a exigência de formalização da ciência moderna, estabelecendo um modelo de decodificação e salvaguarda como ponte para o reconhecimento de direitos e a valorização de tecnologias sociais ancestrais.


\section{Problema de Pesquisa}

O problema central reside na falha de mercado decorrente da assimetria informacional entre os detentores de ativos tradicionais e os agentes econômicos e regulatórios. Embora os sistemas produtivos quilombolas gerem serviços ecossistêmicos e produtos de alta complexidade (ativos VRIN) \cite{Barney1991}, a inexistência de mecanismos formais de mensuração e sinalização (\textit{signaling}) impede a captura desse valor intangível. Sem ferramentas auditáveis capazes de converter conhecimento tácito e difuso em parâmetros verificáveis, esses ativos permanecem subprecificados, vulneráveis à apropriação indevida e excluídos de mercados de maior valor agregado.

Essa lacuna desdobra-se em fricções institucionais e metodológicas que se retroalimentam. A assimetria regulatória e os custos de transação associados aos processos de \textit{compliance} criam barreiras de entrada para comunidades com baixa capacidade técnico-institucional \cite{Agrawal2002}. Paralelamente, verifica-se o esgotamento de métodos lineares de elicitação para capturar fronteiras difusas e dimensões simbólicas do conhecimento tradicional, demandando modelagens computacionais capazes de operar com ambiguidade e incerteza \cite{Zadeh1965}. A ausência de sinais de mercado confiáveis (origem, autenticidade e sustentabilidade) sustenta dinâmicas de seleção adversa, reduzindo incentivos à preservação e à governança territorial \cite{Belletti2015}.


\section{Questões de Pesquisa}

A primeira questão (Q1) investiga em que medida as técnicas de aprendizado de máquina aplicadas a sistemas agrícolas tradicionais apresentam prontidão operacional suficiente para sustentar aplicações regulatórias e de governança territorial, considerando as lacunas na governança de dados, conformidade FAIR e protocolos de validação espacial dos estudos primários. A segunda questão (Q2) indaga se as oito dimensões funcionais de vulnerabilidade biocultural ($V_{1}$--$V_{8}$) apresentam magnitudes de efeito equivalentes ou heterogêneas diante de intervenções de salvaguarda, e em que medida essa heterogeneidade permite hierarquizar as dimensões para subsidiar a ponderação de variáveis no ISB. A terceira questão (Q3) examina como métodos estruturados de adaptação transcultural (WOCAT-SLM) e elicitação (Delphi) contribuem para converter saberes tácitos em variáveis mensuráveis e computacionalmente processáveis, e se o índice de consenso estatístico obtido supera o de métodos de levantamento qualitativo não estruturados. A quarta questão (Q4) examina se o Índice de Salvaguarda Biocultural (ISB), construído pela integração de lógica fuzzy e modelagem psicométrica (TRI) sobre as variáveis elicitadas e ponderadas nas fases precedentes, produz rankings de prioridade de salvaguarda com maior concordância com a percepção comunitária de urgência do que métodos de classificação linear ou ordinal. A quinta questão (Q5) investiga se uma arquitetura computacional de auditoria normativa, integrando inferência probabilística (Rede Bayesiana), verificação semântica (NLP/SBERT) e ancoragem criptográfica (DLT permissionada), é capaz de verificar automaticamente a conformidade dos dossiês de salvaguarda produzidos pelo modelo (dossiês ISB, fichas WOCAT-SLM-QBR, relatórios fuzzy) frente aos requisitos de instâncias regulatórias (IPHAN, INPI, CGen, FAO-GIAHS), com padronização, completude e rastreabilidade superiores às da avaliação manual convencional. A sexta questão (Q6) investiga se a incorporação de usuários finais na construção das regras e interfaces do modelo (\textit{participatory design}) resulta em taxa de adoção continuada da tecnologia de decodificação superior à observada em ferramentas de gestão exógenas.


\section{Hipóteses}

A hipótese H1 postula que a prontidão operacional das técnicas de aprendizado de máquina aplicadas a sistemas agrícolas tradicionais é insuficiente para sustentar aplicações regulatórias e de governança territorial, em razão de lacunas sistêmicas na governança de dados dos estudos primários.

A hipótese H2 postula que as oito dimensões funcionais de vulnerabilidade biocultural: erosão intergeracional e migração ($V_{1}$), complexidade e singularidade biocultural ($V_{2}$), status de documentação ($V_{3}$), vulnerabilidade jurídica e fundiária ($V_{4}$), organização social e governança ($V_{5}$), vitalidade linguística ($V_{6}$), integração ao mercado ($V_{7}$) e exposição climática ($V_{8}$), não apresentam magnitudes de efeito equivalentes diante de intervenções de salvaguarda, de modo que ao menos uma dessas dimensões difere significativamente das demais quanto à sua magnitude de efeito.

A hipótese H3 sustenta que a aplicação combinada de adaptação transcultural do questionário WOCAT-SLM e método Delphi produz variáveis computacionalmente processáveis cujo índice de consenso estatístico é superior ao obtido por métodos de levantamento qualitativo não estruturados.

A hipótese H4 propõe que o Índice de Salvaguarda Biocultural (ISB), construído pela integração de lógica fuzzy e modelagem psicométrica (TRI) sobre variáveis elicitadas e validadas nas fases precedentes, produz rankings de prioridade de salvaguarda com maior concordância com a percepção comunitária de urgência do que métodos de classificação linear ou ordinal.

A hipótese H5 postula que dossiês de salvaguarda auditados por uma arquitetura computacional integrada (Rede Bayesiana para conformidade estrutural, NLP/SBERT para verificação semântica e DLT permissionada para ancoragem criptográfica) apresentam padronização, completude e rastreabilidade documental superiores às da documentação avaliada por procedimentos manuais convencionais.

A hipótese H6 prevê que a incorporação dos usuários finais na etapa de design das regras de inferência (\textit{participatory design}) resulta em taxa de adoção continuada da tecnologia de decodificação superior à observada em ferramentas de gestão exógenas.



\section{Objetivos}
\subsection{Objetivo Geral}

Desenvolver e validar um modelo integrado de decodificação e salvaguarda, articulando aprendizado de máquina, lógica fuzzy e modelagem psicométrica, para documentar e proteger ativos intangíveis em Saberes e Sistemas Agrícolas Tradicionais de comunidades quilombolas do Território de Identidade Semiárido Nordeste II (Bahia).

\subsection{Objetivos Específicos}

O primeiro objetivo específico (OE1) consiste em mapear e avaliar criticamente o estado da arte na aplicação de aprendizado de máquina para monitoramento e avaliação de sistemas agrícolas tradicionais, identificando lacunas metodológicas em validação espacial, explicabilidade e governança de dados que condicionam a prontidão operacional dessas tecnologias. O segundo objetivo (OE2) visa investigar, por meio de meta-análise sob diretrizes PRISMA~2020, como diferentes dimensões de vulnerabilidade biocultural se associam à degradação de saberes e sistemas agrícolas tradicionais, avaliando a variabilidade dessas associações e identificando fatores contextuais que as modulam, de modo a subsidiar a ponderação de variáveis no ISB. O terceiro objetivo (OE3) visa converter saberes tácitos em variáveis computacionalmente processáveis e auditáveis mediante adaptação transcultural do questionário WOCAT-SLM e técnicas de elicitação estruturada (Delphi), aumentando consenso e comparabilidade entre avaliadores. O quarto objetivo (OE4) propõe desenvolver e implementar um motor de inferência baseado em lógica fuzzy e modelagem psicométrica (TRI), integrando as variáveis ponderadas para gerar o Índice de Salvaguarda Biocultural (ISB), classificando saberes segundo urgência de documentação e proteção. O quinto objetivo (OE5) busca avaliar a eficácia probatória da documentação gerada pelo modelo (dossiês ISB, fichas WOCAT-SLM-QBR, relatórios fuzzy) mediante arquitetura computacional de auditoria normativa que integra inferência probabilística (Rede Bayesiana), verificação semântica (NLP/SBERT) e ancoragem criptográfica (DLT permissionada), verificando conformidade normativa e rastreabilidade perante órgãos de salvaguarda e proteção de propriedade intelectual (IPHAN, INPI, CGen, FAO-GIAHS). O sexto objetivo (OE6) consiste em validar a adoção territorial do modelo de salvaguarda mediante design participativo com comunidades quilombolas, avaliando usabilidade, apropriação e sustentabilidade institucional da tecnologia como instrumento de empoderamento e proteção de ativos bioculturais.


\section{Justificativa}

A vulnerabilidade dos Saberes e Sistemas Agrícolas Tradicionais (SSAT) constitui um fenômeno multidimensional cuja escala e velocidade de progressão impõem urgência científica e regulatória. A erosão intergeracional e migração ($V_{1}$), expressa pela aculturação (substituição \textit{in situ} de práticas endógenas por padrões da cultura dominante) e pela emigração de detentores para centros urbanos, cujo efeito combinado reduz a proporção de jovens que dominam rotinas agroecológicas entre coortes consecutivas, compromete a própria via de transmissão oral pela qual os SSAT se reproduzem \cite{ReyesGarcia2013,GomezBaggethun2013}. A perda de detentores ativos desencadeia efeito cascata sobre a complexidade e singularidade biocultural ($V_{2}$), pois a simplificação do repertório de manejo reduz a riqueza de espécies cultivadas e elimina interações solo-planta cuja codificação depende de continuidade geracional, ao passo que o grau de endemismo ecológico e cultural dos agroecossistemas amplifica essa fragilidade porque a perda de um único sistema pode significar a extinção irreversível de genótipos e práticas sem equivalentes em outras paisagens \cite{Jarvis2008,Brush2004,Bussmann2018,CalvetMir2016}. Paralelamente, o status de documentação ($V_{3}$) permanece criticamente deficitário, uma vez que a natureza tácita e oral dos saberes impede a instrução de processos formais de salvaguarda patrimonial (IPHAN), registro de propriedade intelectual coletiva (INPI) \cite{INPI2023}, autorização de acesso ao patrimônio genético (CGen) \cite{brasil2015lei13123} e candidatura a programas internacionais de preservação (FAO-GIAHS). A vulnerabilidade jurídica e fundiária ($V_{4}$), evidenciada pela associação entre ausência de titulação territorial e violência contra comunidades quilombolas \cite{conaq2023racismo}, agrava-se quando inexiste documentação técnica padronizada que comprove ocupação ancestral e gestão sustentável de recursos, fragilizando a posição comunitária em disputas fundiárias e processos de licenciamento \cite{Santilli2009}. A degradação da organização social e governança ($V_{5}$) retroalimenta todas as dimensões precedentes, pois a desarticulação das estruturas coletivas de governança enfraquece os mecanismos endógenos de transmissão, controle e sanção que sustentam a integridade do sistema \cite{Agrawal2002,Folke2016}. A vitalidade linguística ($V_{6}$) opera como veículo de codificação dos saberes, de modo que a substituição de línguas e dialetos locais por idiomas dominantes compromete a nomenclatura etnobotânica e os protocolos orais de manejo, acelerando a perda de precisão na transmissão do conhecimento ecológico local \cite{CamaraLeret2021,Maffi2001}. A integração ao mercado ($V_{7}$) pode atuar como fator de proteção quando viabiliza cadeias de valor que incentivam a manutenção de práticas tradicionais, porém tende a operar como vetor de vulnerabilidade quando a pressão por padronização comercial induz a substituição de cultivares locais por variedades de alto rendimento \cite{Godoy2005,ReyesGarcia2013}. A exposição climática ($V_{8}$) afeta diretamente a viabilidade dos agroecossistemas tradicionais, pois a alteração nos regimes de precipitação, temperatura e eventos extremos desestabiliza os calendários fenológicos ancestrais e os mecanismos adaptativos que sustentam a resiliência socioecológica \cite{Savo2016,Aniah2019}.

A interdependência não linear dessas oito dimensões configura um problema de engenharia de sistemas que métodos setoriais não conseguem resolver. Abordagens unidimensionais, que se concentrem exclusivamente em documentação etnobotânica, em titulação fundiária ou em digitalização de acervos, tratam sintomas isolados sem intervir nos mecanismos de retroalimentação que aceleram a degradação conjunta. A inexistência de ferramentas integradas capazes de quantificar simultaneamente as magnitudes de efeito de cada dimensão $V_{1}$--$V_{8}$, hierarquizá-las por sensibilidade a intervenções de salvaguarda e converter essa síntese em variáveis computacionais operacionais constitui a lacuna metodológica central que esta tese visa preencher.

O modelo integrado proposto responde a essa lacuna ao articular, em cadeia processual, diagnóstico tecnológico, síntese de evidências, elicitação estruturada, modelagem computacional, auditoria documental e validação territorial, convertendo progressivamente complexidade biocultural em evidência auditável. O diagnóstico da prontidão operacional do aprendizado de máquina identifica as lacunas em governança de dados que condicionam a confiabilidade de modelos preditivos aplicados a SAT \cite{Wilkinson2016,Rudin2019}, enquanto a síntese meta-analítica das magnitudes de efeito das dimensões $V_{1}$--$V_{8}$ e seus moderadores contextuais produz a hierarquização empírica que fundamenta a ponderação de variáveis no Índice de Salvaguarda Biocultural. A adaptação transcultural do WOCAT-SLM e a elicitação Delphi superam as limitações de métodos qualitativos não estruturados ao converter saberes tácitos em variáveis linguísticas validadas por consenso estatístico \cite{WOCAT2023,Linstone1975}, e a modelagem fuzzy-psicométrica articula a Teoria dos Conjuntos Difusos \cite{Zadeh1965} e a Teoria de Resposta ao Item para representar fronteiras nebulosas e ambiguidade cultural, gerando o ISB como escore contínuo de prioridade com aderência à percepção comunitária de urgência. A auditoria computacional assegura que os dossiês produzidos atendam aos requisitos probatórios de instâncias regulatórias (IPHAN, INPI, CGen, FAO-GIAHS), e o design participativo incorpora usuários finais na co-criação de regras e interfaces, condicionando a adoção continuada da tecnologia e evitando a rejeição de ferramentas exógenas documentada na literatura \cite{Schuler1993}.

Além do mérito científico, a arquitetura modular da tese é desenhada para capilaridade acadêmica e formação de competências \cite{Creswell2018}. Cada módulo produz um artefato metodológico autônomo (protocolo PRISMA, modelo meta-analítico, instrumento WOCAT-SLM-QBR adaptado, motor fuzzy-TRI, dossiê auditável, kit de design participativo) capaz de ser replicado em outros territórios e contextos de patrimônio biocultural, ao mesmo tempo em que a integração do conjunto gera um arcabouço de suporte à decisão com potencial de impacto territorial e institucional.

\section{Delineamento do estudo}

Este estudo configura-se como pesquisa aplicada de métodos mistos, com delineamento sequencial e abordagem orientada à construção e validação de um modelo integrado de decodificação e salvaguarda. O delineamento metodológico está estruturado em seis fases integradas. A primeira fase compreende a revisão de escopo e diagnóstico tecnológico, conduzida mediante revisão sistemática (PRISMA-ScR) da literatura sobre aprendizado de máquina aplicado a sistemas agrícolas tradicionais, com meta-análise de acurácia, avaliação de conformidade FAIR e identificação de lacunas em validação espacial e explicabilidade. A segunda fase consiste na meta-análise quantitativa sob diretrizes PRISMA~2020, investigando como diferentes dimensões de vulnerabilidade biocultural se associam à degradação de saberes e sistemas agrícolas tradicionais, avaliando a variabilidade dessas associações e identificando fatores contextuais que as modulam, para subsidiar a ponderação de variáveis no ISB. A terceira fase, dedicada à elicitação e estruturação de variáveis, envolve a adaptação transcultural do questionário WOCAT-SLM e a aplicação de procedimentos Delphi para converter saberes tácitos em variáveis linguísticas auditáveis e critérios mensuráveis. A quarta fase abrange a modelagem fuzzy e a construção do ISB, incluindo o desenho do motor de inferência fuzzy (funções de pertinência, regras, defuzzificação) e o cálculo do Índice de Salvaguarda Biocultural, que classifica saberes segundo urgência de proteção. A quinta fase contempla a auditoria computacional da documentação gerada pelo modelo mediante arquitetura que integra três módulos, Rede Bayesiana para verificação de completude estrutural, NLP/SBERT para aderência semântica dos campos textuais aos descritores normativos e DLT permissionada para ancoragem criptográfica e conformidade FAIR, verificando padronização, completude, rastreabilidade e conformidade normativa perante as quatro instâncias regulatórias (IPHAN, INPI, CGen, FAO-GIAHS). A sexta fase consiste no design participativo e validação territorial, promovendo a co-criação das regras e interfaces do modelo de salvaguarda com usuários finais, avaliação de usabilidade, apropriação e adoção, e consolidação do ecossistema de pesquisa aplicada e formação.



